\newif\ifglobal

%%%%%%%%%%%%%%%%%%%%%%%
%%% Compile to view %%%
%%%%%%%%%%%%%%%%%%%%%%%

%%%%%%%%%%%%%%%%%%%%%%%%%%%%%%%%%%%%%%%%%%%%%%%%%%%
%%% Readme:                                     %%%
%%% https://www.overleaf.com/read/bwdjvfjksvjy/ %%%
%%%%%%%%%%%%%%%%%%%%%%%%%%%%%%%%%%%%%%%%%%%%%%%%%%%

% >>> M?: marks a module setting number or important notice

% >>> M4: Set {./relative-path} to external files for this module <<<
\def \modulefiles {./16-SHA}
% To add an external file, use:
% (Replace '---' with actual filename)
% '\input{\modulefiles/tables/---.tex}' for tables
% '\includegraphics{\modulefiles/figures/---}' for figures
% '\subfile{\modulefiles/---}' for register files


% >>> M1: This module is in global project? <<<
% 'YES' - uncomment; 'NO' - comment out
\globaltrue

\ifglobal
    \documentclass[main\_\_CN.tex]{subfiles}

\else
    % Font "FandolSong-Regular" used by ctex does not have GSUB table.
    % It causes warnings that do not affect compilation and can be ignored.
    \PassOptionsToPackage{quiet}{fontspec}

    \documentclass[openany, 10pt]{book}

    % >>> M2: In ./00-shared/preamble-trm-module, also find
    %         settings for visibility of todo notes, watermark <<<
    \usepackage{preamble-trm-module}

    % Enable Chinese typesetting
    \usepackage{ctex}

    % >>> M3: Temporary labels in Chinese
    %         you can add your own labels there <<<
    \myexternaldocument{./00-shared/temp-labels-cn}

\fi %\ifglobal


\setcounter{secnumdepth}{4}


% >>> M5: If this module needs any updates in
% './00-shared/preamble-trm-module.sty'
% (include additional package, etc.), contact your module PM
% or create the file './00-shared/changelog.tex' make the
% necessary updates yourself and document them in this file <<<

% Make text left-aligned and allow latex to break pages naturally, comment out for full justification
\raggedright
\raggedbottom


\begin{document}


% Sets language into which pre-defined bits of text are translated
\selectlanguage{chinese}
\setchineseformatting


% Do not delete this block
\ifglobal
% Add the following front matter if
% this module is outside of global project
\else
    \InputIfFileExists{readme}{}{}
    \listoftodos
    {\let\clearpage\relax \listofdonetodos}
    \tableofcontents
    %\listoftables
    %\listoffigures
\fi


%%%%%%%%%%%%%%%%%%%%%%%%%
%%% TRM Module Begins %%%
%%%%%%%%%%%%%%%%%%%%%%%%%

\hypertarget{sha}{}
\chapter{SHA 加速器 (SHA)}
\label{mod:sha}


\section{概述}
相比于在软件中单独运行 SHA （安全哈希算法），SHA 加速器能够快速提升 SHA 的运行速度。SHA 加速器支持四种标准 FIPS PUB 180-4 运算：SHA-1、SHA-256、SHA-384 和 SHA-512。


\section{主要特性}

\begin{itemize}
    \item 支持SHA-1 运算
    \item 支持SHA-256 运算
    \item 支持SHA-384 运算
    \item 支持SHA-512 运算
\end{itemize}


\section{功能描述}



\subsection{填充解析信息}

SHA 加速器每次只能接受一个信息块。软件将整个信息按照标准“FIPS PUB 180-4”中“5.2 Parsing the Message”的要求划分为一个一个的信息块，每次只将一个信息块写入寄存器 SHA\_TEXT\_\textcolor{red}{\it{n}}\_REG。如果是 SHA-1、SHA-256 运算，则软件每次将 512 bit 的块写入寄存器 SHA\_TEXT\_0\_REG\verb+ ~ +SHA\_TEXT\_15\_REG。
如果是 SHA-384、SHA-512 运算，则软件每次将 1024 bit 的块写入寄存器 SHA\_TEXT\_0\_REG\verb+ ~ +SHA\_TEXT\_31\_REG。

SHA 加速器不能自动完成标准“FIPS PUB 180-4”中“5.1 Padding the Message”所要求的填充操作；在将信息输入到加速器之前，需由软件来完成填充操作。

标准“FIPS PUB 180-4”中“2.2.1 Parameters”描述“$M_{0}^{(i)}$ is the left-most word of message block i”，此 word 存放在寄存器 SHA\_TEXT\_0\_REG 中。以此类推寄存器 SHA\_TEXT\_1\_REG 中存放的是信息中某个块从左向右的第二个 word……


\subsection{信息摘要}

哈希运算完成之后，信息摘要被 SHA 加速器更新入寄存器 SHA\_TEXT\_\textcolor{red}{\it{n}}\_REG。
如果是 SHA-1 运算，则 160 bit 信息摘要存放在寄存器 SHA\_TEXT\_0\_REG\verb+ ~ +SHA\_TEXT\_4\_REG。
如果是 SHA-256 运算，则 256 bit 信息摘要存放在寄存器 SHA\_TEXT\_0\_REG\verb+ ~ +SHA\_TEXT\_7\_REG。
如果是 SHA-384 运算，则 384 bit 信息摘要存放在寄存器 SHA\_TEXT\_0\_REG\verb+ ~ +SHA\_TEXT\_11\_REG。
如果是 SHA-512 运算，则 512 bit 信息摘要存放在寄存器 SHA\_TEXT\_0\_REG\verb+ ~ +SHA\_TEXT\_15\_REG。

标准“FIPS PUB 180-4”中“2.2.1 Parameters”描述“$H^{(N)}$ is the final hash value and is used to determine the message digest”、“$H_{0}^{(i)}$ is the left-most word of hash value i”。
则信息摘要中最左边 word  $H_{0}^{(N)}$ 在寄存器 SHA\_TEXT\_0\_REG 中。以此类推信息摘要中从左至右的第二个 word $H_{1}^{(N)}$ 在寄存器 SHA\_TEXT\_1\_REG 中……


\subsection{哈希运算}

运算 SHA-1、SHA-256、SHA-384 和 SHA-512 各有一组控制寄存器；不同的哈希运算使用不同的控制寄存器。
运算 SHA-1 使用寄存器 SHA\_SHA1\_START\_REG、SHA\_SHA1\_CONTINUE\_REG、SHA\_SHA1\_LOAD\_REG 和 SHA\_SHA1\_BUSY\_REG。
运算 SHA-256 使用寄存器 SHA\_SHA256\_START\_REG、SHA\_SHA256\_CONTINUE\_REG、SHA\_SHA256\_LOAD\_REG 和 SHA\_SHA256\_BUSY\_REG。
运算 SHA-384 使用寄存器 SHA\_SHA384\_START\_REG、SHA\_SHA384\_CONTINUE\_REG、SHA\_SHA384\_LOAD\_REG 和 SHA\_SHA384\_BUSY\_REG。
运算 SHA-512 使用寄存器 SHA\_SHA512\_START\_REG、SHA\_SHA512\_CONTINUE\_REG、SHA\_SHA512\_LOAD\_REG 和\\ SHA\_SHA512\_BUSY\_REG。
具体操作流程如下。

\begin{enumerate}
    \item 操作第一个信息块
    \begin{enumerate}
        \item 用第一个信息块初始化寄存器 SHA\_TEXT\_\textcolor{red}{\it{n}}\_REG。
		\item 对寄存器 SHA\_\textcolor{red}{\it{X}}\_START\_REG 写入 1。
		\item 轮询寄存器 SHA\_\textcolor{red}{\it{X}}\_BUSY\_REG，直到从寄存器 SHA\_\textcolor{red}{\it{X}}\_BUSY\_REG 读出 0。
    \end{enumerate}

    \item 循环操作信息的后续每一个块
    \begin{enumerate}
        \item 用后续的信息块初始化寄存器 SHA\_TEXT\_\textcolor{red}{\it{n}}\_REG。
		\item 对寄存器 SHA\_\textcolor{red}{\it{X}}\_CONTINUE\_REG 写入 1。
		\item 轮询寄存器 SHA\_\textcolor{red}{\it{X}}\_BUSY\_REG，直到从寄存器 SHA\_\textcolor{red}{\it{X}}\_BUSY\_REG 读出 0。
    \end{enumerate}

    \item 获取信息摘要
    \begin{enumerate}
        \item 对寄存器 SHA\_\textcolor{red}{\it{X}}\_LOAD\_REG 写入 1。
		\item 轮询寄存器 SHA\_\textcolor{red}{\it{X}}\_BUSY\_REG，直到从寄存器 SHA\_\textcolor{red}{\it{X}}\_BUSY\_REG 读出 0。
		\item 从寄存器 SHA\_TEXT\_\textcolor{red}{\it{n}}\_REG 取出信息摘要。
	\end{enumerate}
\end{enumerate}

\subsection{运行效率}

SHA 加速器需要 60 至 100 个时钟周期来处理一个信息块以及 8 至 20 个时钟周期来计算最后的信息摘要。

\hypertarget{sha-reg-summ}{}
\section{寄存器列表}

本小节的所有地址均为相对于 SHA 加速器基地址的地址偏移量（相对地址），具体基地址请见章节 \ref{mod:sysmem} \textit{\nameref{mod:sysmem}} 中的表 \ref{table:Peripheral_Address_Mapping}。

请查看章节 \hyperref[glossary-access-types]{\textit{寄存器的访问类型}}，了解“访问”列缩写的含义。

\subfile{\modulefiles/sha_reg_summary_cn}

\newpage

\section{寄存器}

本小节的所有地址均为相对于 SHA 加速器基地址的地址偏移量（相对地址），具体基地址请见章节 \ref{mod:sysmem} \textit{\nameref{mod:sysmem}} 中的表 \ref{table:Peripheral_Address_Mapping}。

关于保留 (reserved) 域的处理，请查看章节 \hyperref[programming-reserved-field]{\textit{如何配置寄存器的保留域}}。

\subfile{\modulefiles/sha_reg_cn}

\end{document}